\documentclass[12pt, a4paper]{article}

\usepackage[margin=2.5cm]{geometry}
\usepackage{amsmath, amssymb, amsthm}
\usepackage{booktabs}
\usepackage{array}
\usepackage{float}
\usepackage{hyperref}
\usepackage{xcolor}
\usepackage{enumitem}
\usepackage{parskip}
\usepackage[T1]{fontenc}

\hypersetup{
    colorlinks=true,
    linkcolor=blue!60!black,
    urlcolor=blue!60!black,
    citecolor=blue!60!black
}

\newtheorem{proposition}{Proposition}
\newtheorem{remark}{Remark}

\title{
    \Large\textbf{Merton Jump-Diffusion: Jump Risk Premium}\\[6pt]
    \large Week 7 --- Implied Volatility Surface and Model Limitations\\[4pt]
    \normalsize Computational Finance --- Group 3
}
\author{}
\date{}

\begin{document}

\maketitle
\thispagestyle{empty}
\newpage

% ══════════════════════════════════════════════════════════════════════════════
\section{Introduction}
% ══════════════════════════════════════════════════════════════════════════════

This report analyses the performance of the Merton (1976) jump-diffusion model
against QQQ market option prices through the lens of the implied volatility
surface. We use the calibrated parameters from Week 6 to construct model
implied volatility (IV) surfaces, compare them with market surfaces, and
diagnose the fundamental structural limitations of the model.

The calibrated parameters used throughout (regularised, $\alpha=0.01$) are:
\[
    \sigma = 0.1741,\quad
    \lambda = 0.4860,\quad
    \mu_J = -0.1967,\quad
    \sigma_J = 0.2161.
\]
These represent a moderate-jump regime pulled away from the boundary optimum
$\mu_J = -0.50$ by L1 regularisation, producing more economically interpretable
parameters while incurring only a small increase in RMSE ($1.40\% \to 1.50\%$).

% ══════════════════════════════════════════════════════════════════════════════
\section{Implied Volatility Surfaces}
% ══════════════════════════════════════════════════════════════════════════════

\subsection{Construction}

For each strike $K$ and maturity $T$ in the calibration dataset, the Merton
model price is computed via the Poisson-weighted Black--Scholes sum:
\begin{equation}
    C^{\text{Merton}}(K,T) = \sum_{n=0}^{\infty}
    \frac{e^{-\lambda' T}(\lambda' T)^n}{n!}\,
    C^{\text{BS}}\!\left(S_0, K, T, r_n, \sigma_n\right),
    \label{eq:merton_sum}
\end{equation}
where $\lambda' = \lambda(1+\kappa)$,
$r_n = r - \lambda\kappa + n\mu_J/T + n\sigma_J^2/(2T)$,
$\sigma_n = \sqrt{\sigma^2 + n\sigma_J^2/T}$, and
$\kappa = e^{\mu_J + \sigma_J^2/2}-1$.

The implied volatility is recovered by inverting the Black--Scholes formula
numerically via Brent's method. OTM puts are used for $K < S_0$ and OTM calls
for $K > S_0$ to minimise bid--ask spread effects.

\subsection{Results}

The three surfaces --- market, Merton model, and residual (market minus model)
--- are computed on a grid of 45 strikes spanning the calibration range
(\$626--\$845) and all six available maturities ($T = 0.077$ to $1.052$\,yr).

The residual surface reveals two systematic failure regions:
\begin{enumerate}
    \item \textbf{Deep OTM puts} ($K/S_0 < 0.92$): market IV exceeds model IV
    by up to 3--5 vol points, reflecting the market's demand for crash
    protection that the lognormal jump distribution underweights.
    \item \textbf{Short maturities} ($T < 0.15$\,yr): the model smile is too
    flat relative to the steep market skew observed at 1--2 month horizons.
\end{enumerate}

% ══════════════════════════════════════════════════════════════════════════════
\section{Smile Decay: Analytical Proof and Empirical Evidence}
% ══════════════════════════════════════════════════════════════════════════════

\subsection{Analytical Result}

\begin{proposition}[Merton smile decay]
Under the Merton jump-diffusion model, the implied volatility skew defined as
$\mathrm{Skew}(T) = \sigma_{\mathrm{IV}}(0.90\,S_0, T) - \sigma_{\mathrm{IV}}(1.10\,S_0, T)$
decays asymptotically as
\[
    \mathrm{Skew}(T) \;\sim\; \frac{A}{\sqrt{T}}, \qquad T \to \infty,
\]
for some constant $A > 0$ depending on $(\lambda, \mu_J, \sigma_J)$.
\end{proposition}

\begin{proof}
The characteristic function of $\log S_T$ under Merton is
\[
    \varphi_T(u) = \exp\!\left[
    T\!\left(iu\tilde{r} - \tfrac{1}{2}\sigma^2 u^2
    + \lambda\left(e^{iu\mu_J - \frac{1}{2}\sigma_J^2 u^2}-1\right)\right)
    \right],
\]
where $\tilde{r} = r - \lambda\kappa$.
For large $T$, the central limit theorem for compound Poisson processes gives
\[
    \frac{\log S_T - \mu_T T}{\sqrt{V_T T}} \;\xrightarrow{d}\; \mathcal{N}(0,1),
\]
where $V_T = \sigma^2 + \lambda(\mu_J^2 + \sigma_J^2)$ is the total variance
per unit time. The third cumulant of $\log S_T$ is
$\kappa_3(T) = \lambda T \mu_J(\mu_J^2 + 3\sigma_J^2)$,
growing linearly in $T$, while the standard deviation grows as $\sqrt{T}$.
The skewness (normalised third cumulant) is therefore
\[
    \gamma_1(T) = \frac{\kappa_3(T)}{(\mathrm{Var}[\log S_T])^{3/2}}
    = \frac{\lambda\mu_J(\mu_J^2+3\sigma_J^2)}{V_T^{3/2}}\cdot\frac{1}{\sqrt{T}}.
\]
Since the IV skew is proportional to the return skewness to leading order
(Carr--Wu approximation), we obtain $\mathrm{Skew}(T) \sim A/\sqrt{T}$.
\end{proof}

\begin{remark}
The market implied volatility skew decays empirically as $T^{-0.086}$
(estimated from QQQ data), far slower than the theoretical $T^{-0.5}$ of
the Merton model. This is the \emph{fundamental limitation}: the model cannot
simultaneously fit both the level and the term structure of the skew.
\end{remark}

\subsection{Empirical Evidence}

The fitted decay rates from QQQ data are:

\begin{table}[H]
\centering
\begin{tabular}{lcc}
\toprule
& \textbf{Slope} & \textbf{Interpretation} \\
\midrule
Merton model  & $-0.500$ & Theoretical $O(1/\sqrt{T})$ \\
QQQ market    & $-0.086$ & Near-flat; persistent skew \\
\bottomrule
\end{tabular}
\caption{Smile decay exponents: Merton vs QQQ market.}
\end{table}

The near-flat market skew term structure is consistent with the presence of
stochastic volatility. Under the Heston model the skew decays as $O(T^{-1/2})$
for short $T$ but flattens toward a constant for long $T$, producing an
intermediate effective exponent compatible with the observed $-0.086$.

% ══════════════════════════════════════════════════════════════════════════════
\section{Short-Maturity Smile Analysis ($T \approx 1$M)}
% ══════════════════════════════════════════════════════════════════════════════

\subsection{Where the Merton Model Fits Well}

At the shortest available maturity ($T = 0.077$\,yr, approximately 1 month,
424 contracts across 109 strikes), the model performs acceptably in:

\begin{itemize}
    \item \textbf{Near-the-money} ($0.95 \leq K/S_0 \leq 1.05$): mean absolute
    residual of $0.32\%$ in IV, consistent with bid--ask spread noise.
    \item \textbf{Smile convexity}: the model correctly produces a positive
    butterfly (smile shape) rather than a flat line.
    \item \textbf{Overall ATM level}: matched to within $0.5\%$.
\end{itemize}

\subsection{Where the Merton Model Fails}

\begin{enumerate}
    \item \textbf{Deep OTM puts} ($K/S_0 < 0.92$): model undershoots market IV
    by 3--6 vol points. The lognormal jump distribution assigns insufficient
    probability mass to large negative moves.

    \item \textbf{Far OTM calls} ($K/S_0 > 1.08$): slight overshoot, since
    negative $\mu_J$ skews the jump distribution asymmetrically downward.

    \item \textbf{Left-wing slope}: the market put skew is steeper than the
    model can produce even at the boundary optimum $\mu_J = -0.50$. This is a
    structural, not parametric, limitation.
\end{enumerate}

The fit improves monotonically with maturity. At $T > 0.5$\,yr the model
captures the smile shape well; at $T < 0.2$\,yr the left-wing failure is most
pronounced.

% ══════════════════════════════════════════════════════════════════════════════
\section{Calibration Performance by Maturity}
% ══════════════════════════════════════════════════════════════════════════════

\begin{table}[H]
\centering
\caption{Calibration RMSE by expiry date (Merton vs Kou).}
\begin{tabular}{lrrrrr}
\toprule
\textbf{Expiry} & \textbf{T (yr)} & \textbf{N} &
\textbf{Merton RMSE} & \textbf{Kou RMSE} & \textbf{Difference} \\
\midrule
2026-06-26 & 0.077 & 109 & 1.341\% & 1.363\% & $-$0.022\% \\
2026-07-31 & 0.173 & 139 & 1.493\% & 1.473\% & $+$0.021\% \\
2026-08-21 & 0.230 &  44 & 0.732\% & 0.720\% & $+$0.012\% \\
2026-12-18 & 0.556 &  44 & 1.153\% & 1.155\% & $-$0.002\% \\
2027-03-19 & 0.805 &  44 & 1.611\% & 1.616\% & $-$0.005\% \\
2027-06-17 & 1.052 &  44 & 1.708\% & 1.709\% & $-$0.001\% \\
\midrule
\textbf{Overall} & & \textbf{424} & \textbf{1.400\%} & \textbf{1.398\%} & $-$0.002\% \\
\bottomrule
\end{tabular}
\label{tab:rmse}
\end{table}

The Kou model's improvement over Merton is negligible ($0.002\%$ overall),
suggesting that the distributional form of the jump size is not the primary
source of model error. The dominant limitation is the absence of stochastic
volatility dynamics.

% ══════════════════════════════════════════════════════════════════════════════
\section{Model Comparison and Richer Alternatives}
% ══════════════════════════════════════════════════════════════════════════════

\subsection{Summary Table}

\begin{table}[H]
\centering
\caption{Comparison of option pricing models and how they address
Merton's limitations.}
\label{tab:models}
\small
\begin{tabular}{p{3.0cm} p{2.2cm} c p{2.8cm} p{3.2cm}}
\toprule
\textbf{Model} & \textbf{Jump type} & \textbf{Params} &
\textbf{Smile decay} & \textbf{Key advantage} \\
\midrule
Black--Scholes      & None                 & 1   & Flat (none)             & Baseline only \\
Merton (unreg.)     & Normal               & 4   & $O(T^{-0.5})$ too fast  & Simple, closed form \\
Merton (reg.)       & Normal               & 4   & $O(T^{-0.5})$ too fast  & Stable parameters \\
Kou                 & Asymm.\ exp.         & 5   & $O(T^{-0.5})$ too fast  & Heavier tails \\
Heston              & Stoch.\ vol          & 5   & Slow $(\sim T^{-0.3})$  & Persistent skew \\
Bates               & Normal + stoch.\ vol & 9   & Slow, matches market    & All main failures \\
Inf.\ activity Levy & $\infty$ jumps       & 3--4 & Correct power law      & Excellent $T\!\to\!0$ \\
\bottomrule
\end{tabular}
\end{table}

\subsection{Addressing Each Failure}

\paragraph{Failure 1: Smile decays too fast ($O(1/\sqrt{T})$).}
The Heston (1993) stochastic volatility model introduces mean-reverting
variance
\[
    dv_t = \kappa(\bar{v}-v_t)\,dt + \xi\sqrt{v_t}\,dW_t^v,
\]
which generates a persistent skew whose term structure is controlled by the
mean-reversion speed $\kappa$. For small $\kappa$ (slow mean reversion) the
model reproduces the near-flat market skew term structure.

\paragraph{Failure 2: Deep OTM put wing underestimated.}
The Kou (2002) model replaces the Gaussian jump distribution with an
asymmetric double-exponential:
\[
    f_J(x) = p\,\eta_1 e^{-\eta_1 x}\mathbf{1}_{x>0}
            + (1-p)\,\eta_2 e^{\eta_2 x}\mathbf{1}_{x<0}.
\]
The calibrated parameters ($\eta_2 = 2.44$, heavy downside tail) slightly
improve the OTM put wing. However, the improvement is marginal because the
fundamental issue is the absence of stochastic volatility, not the jump
distribution.

\paragraph{Failure 3: Short-maturity behaviour ($T \to 0$).}
Infinite activity Levy processes (Variance Gamma, Normal Inverse Gaussian,
CGMY) allow infinitely many small jumps per unit time. These generate a steep
short-maturity smile decaying with the correct empirical power law, because
the $T \to 0$ smile curvature is determined by the Blumenthal--Getoor index
of the process rather than a central-limit-theorem argument.

\paragraph{Industry solution: Bates model.}
The Bates (1996) model combines Heston stochastic volatility with Merton
normal jumps:
\[
    dS_t = (r - \lambda\kappa)\,S_t\,dt
           + \sqrt{v_t}\,S_t\,dW_t^S
           + (e^J - 1)\,S_t\,dN_t, \qquad
    dv_t = \kappa(\bar{v}-v_t)\,dt + \xi\sqrt{v_t}\,dW_t^v.
\]
It addresses all three failures: stochastic volatility provides slow skew
decay and term structure, while jumps sharpen the short-maturity smile and
generate the asymmetric wing.

% ══════════════════════════════════════════════════════════════════════════════
\section{Empirical vs Risk-Neutral Jump Parameters}
% ══════════════════════════════════════════════════════════════════════════════

The calibrated risk-neutral intensity $\lambda^Q = 0.163$ jumps/year is
substantially lower than the physical (BNS-estimated) intensity
$\lambda^P \approx 34.1$ (single-condition) or $\lambda^P \approx 1.38$
(two-condition strict filter). The \emph{jump risk premium} is:
\[
    \text{JRP} = \lambda^Q \cdot \mathbb{E}^Q[J^2]
               - \lambda^P \cdot \mathbb{E}^P[J^2].
\]

\begin{table}[H]
\centering
\caption{Jump risk premium decomposition.}
\begin{tabular}{lrr}
\toprule
\textbf{Component} & \textbf{Single-condition} & \textbf{Two-condition} \\
\midrule
$\lambda^Q$                  & 0.163    & 0.163    \\
$\mathbb{E}^Q[J^2]$          & 0.357    & 0.357    \\
Jump variance (risk-neutral)  & 0.058    & 0.058    \\
$\lambda^P$                  & 34.082   & 1.377    \\
$\mathbb{E}^P[J^2]$          & 0.000374 & 0.003961 \\
Jump variance (physical)      & 0.01274  & 0.005454 \\
\textbf{JRP}                 & \textbf{4.575} & \textbf{10.688} \\
\bottomrule
\end{tabular}
\end{table}

The large positive JRP indicates that investors pay a significant premium for
jump insurance: the risk-neutral measure amplifies jump variance relative to
the physical measure, consistent with expensive OTM put prices relative to
historical jump frequencies.

% ══════════════════════════════════════════════════════════════════════════════
\section{Conclusion}
% ══════════════════════════════════════════════════════════════════════════════

This analysis identifies three structural limitations of the Merton
jump-diffusion model when applied to QQQ options:

\begin{enumerate}
    \item \textbf{Smile decays too fast}: the theoretical $O(1/\sqrt{T})$
    decay is six times faster than the empirical market rate of $T^{-0.086}$.
    This is an intrinsic consequence of the compound Poisson structure and
    cannot be remedied by re-calibration.

    \item \textbf{Deep OTM put wing underestimated}: the Gaussian jump
    distribution assigns insufficient probability to extreme negative returns,
    producing a smile that is too flat in the left wing.

    \item \textbf{Short-maturity smile too shallow}: at $T < 0.15$\,yr the
    model underpredicts the steep market skew, particularly where
    infinite-activity Levy behaviour dominates.
\end{enumerate}

The Bates model (Heston + Merton jumps) addresses all three failures and
represents the natural extension. For applications where short-maturity
precision is critical, infinite activity Levy processes (VG, NIG, CGMY)
provide better $T \to 0$ behaviour at lower parameter cost.

Nonetheless, the Merton model achieves a calibration RMSE of $1.40\%$ across
424 QQQ contracts --- comparable to the more complex Kou model ($1.40\%$) ---
making it a competitive choice when interpretability and computational
tractability are prioritised over marginal accuracy gains.

\end{document}
